\documentclass[12pt,a4paper]{exam}
\usepackage[T1]{fontenc}
\usepackage{longtable}
\usepackage{graphicx}
\usepackage{hyperref}

%\urlstyle{same}  % don't use monospace font for urls
\setlength{\parindent}{0pt}
\setlength{\parskip}{6pt plus 2pt minus 1pt}

\usepackage{geometry}
\usepackage{multicol}
\usepackage{relsize}
\usepackage{pgf}

\geometry{
  %includeheadfoot,
  margin=2.54cm
}
\title{Exercises on Design Patterns II}
\author{Even more basics}
\date{2018}

\fvset{tabsize=2}
\fvset{fontsize=\relsize{-1}}
%\fvset{frame=single}

%http://hp.fciencias.unam.mx/~alg/bd/patrones.pdf

\newcommand{\code}{/Users/keith/Documents/Courses/sdp/repo/worksheets/dp-two/kotlin/src}
%\newcommand{\codesol}{/Users/keith/Courses/sdp/SDP-2017/exercises/week08sol/src/main/scala}

\begin{document}
\maketitle

In these exercises we will be examining the following design patterns:
\begin{itemize}
\item Strategy,
\item Abstract Factory,
\item Builder,
\item Facade,
\item Bridge, and
\item Composite.
\end{itemize}
The source code snippets listed below are written using the Kotlin language; 
similar versions in Scala can be found on the source 
code repository under the \verb!worksheets/dp-one! folder. 

\section*{Short form questions}

\begin{questions}

\question
\begin{parts}
\part
Briefly describe the Strategy Design Pattern?

\begin{solution}
The Strategy Design Pattern defines a family of algorithms, 
encapsulating each one, and making them interchangeable. 
Strategy lets the algorithm vary independently from the clients that use it	
\end{solution}

\part
When is it appropriate to use the \textsc{Strategy Design Pattern}?

\begin{solution}
Use the Strategy pattern when:
\begin{itemize}
\item
Many related classes differ only in their behaviour. 
Strategies provide a way to configure a class with one of many behaviours.
\item
You need different variants of an algorithm. For example, you might define algorithms 
reflecting different space/time trade-offs. Strategies can be used when these variants 
are implemented as a class hierarchy of algorithms.
\item
An algorithm uses data that clients shouldn't know about. 
Use the Strategy pattern to avoid exposing complex, algorithm- specific data structures.
\item
A class defines many behaviours, and these appear as multiple conditional statements in 
its operations. Instead of many conditionals, move related conditional branches into 
their own \verb!Strategy! class.	
\end{itemize}
\end{solution}

\end{parts}

\question
When is it appropriate to use the \textsc{Abstract Factory} design pattern?

\begin{solution}
Use the \textsc{Abstract Factory} pattern when
\begin{itemize}
\item 
A system should be independent of how its products are created, composed, and represented.
\item 
A system should be configured with one of multiple families of products.
\item
A family of related product objects is designed to be used together, 
and you need to enforce this constraint.
\item
You want to provide a class library of products, and you want to reveal just their interfaces, 
not their implementations.
\end{itemize}	
\end{solution}

\question % Builder Design Pattern 

\noindent ``In general, the details of object construction, such as instantiating and initialising
the components that comprise the object, are kept within the object, often as part of 
its constructor.''

Comment on this statement with reference to \emph{modularity} and 
\emph{construction bloat}.

\begin{solution}
This type of design closely ties the object construction process with the 
components that make up the object. This approach is suitable as long as the 
object under construction is simple and the object construction process is 
definite and always produces the same representation of the object.

However, this design may not be effective when the object being created is 
complex and the series of steps constituting the object creation process can be 
implemented in different ways, thus producing different representations of the object. 
Because the different implementations of the construction process are all kept within the 
object, the object can become bulky (construction bloat) and less modular. Subsequently, 
adding a new implementation or making changes to an existing implementation 
requires changes to the existing code.
\end{solution}

\question % Facade Design Pattern 
\begin{parts}
\part What is the Facade Pattern? 	
\begin{solution}
The Facade Pattern makes a complex interface easier to use, using a Facade class. 
The Facade Pattern provides a unified interface to a set of interface in a subsystem. 
Facade defines a higher-level interface that makes the subsystem easier to use.	
\end{solution}
\part When, and why, would you use the Facade Pattern?
\begin{solution}
\begin{itemize}
\item
You want to provide a simple interface to a complex subsystem. Subsystems often get more 
complex as they evolve. Most patterns, when applied, result in more and smaller classes. 
This makes the subsystem more reusable and easier to customise, but it also becomes harder 
to use for clients that donâ't need to customise it. A facade can provide a simple default 
view of the subsystem that is good enough for most clients. Only clients needing 
more customisability will need to look beyond the facade.
\item 
There are many dependencies between clients and the implementation classes of an abstraction. 
Introduce a facade to decouple the subsystem from clients and other subsystems, 
thereby promoting subsystem independence and portability.
\item
You can layer your subsystems. Use a facade to define an entry point to each subsystem level. 
If subsystems are dependent, then you can simplify the dependencies between them by making 
them communicate with each other solely through their facades.
\end{itemize}	
\end{solution}
\end{parts}

\question
When should one make use of the \textsc{Bridge Design Pattern}?

\begin{solution}
You should use the Bridge Pattern when:
\begin{itemize}
\item
You want to avoid a permanent binding between an abstraction and its implementation. 
This might be the case, for example, when the implementation must be selected or switched 
at run-time.
\item
Both the abstractions and their implementations should be extensible by sub-classing. 
In this case, the Bridge pattern lets you combine the different abstractions and 
implementations and extend them independently.
\item
Changes in the implementation of an abstraction should have no impact on clients; that is, 
their code should not have to be recompiled.
\item
You want to share an implementation among multiple objects (perhaps using reference counting),
and this fact should be hidden from the client.
\end{itemize}
\end{solution}

\question % Composite Design Pattern
\begin{parts}
\part
What is the Composite Pattern?

\begin{solution}
The formal definition of the Composite Pattern says that it allows you to compose 
objects into tree structures to represent part- whole hierarchies. Composite lets clients to treat 
individual objects and compositions of objects uniformly.	
\end{solution}

\part
Under what conditions would you use a Composite Design Pattern?

\begin{solution}
\begin{itemize}
\item When you want to represent part-whole hierarchies of objects.
\item When you want clients to be able to ignore the difference between compositions 
of objects and individual objects. Clients will treat all objects in the composite 
structure uniformly.
\end{itemize}	
\end{solution}

\part
What are the four participants of the Composite Design Pattern?

\begin{solution}
\begin{itemize}
	\item	Component 
	\item	Leaf
	\item 	Composite 
	\item 	Client
\end{itemize}
\end{solution}
\end{parts}

\end{questions}

\section*{Long form questions}

\begin{questions}

\question % Strategy Design Pattern 
Create a text formatter for a text editor. 
A text editor can have different text formatters to format text. 
We can create different text formatters and then pass the required one to the text editor, 
so that the editor will able to format the text as required.

The text editor will hold a reference to a common interface for the text formatter 
and the editorâ's job will be to pass the text to the formatter to format the text.
You are required to implement this outline using the \textsc{Strategy Design Pattern} 
which will make the code flexible and maintainable. 

Below is the \verb!TextFormatter! interface which is implemented by all the concrete formatters.

\inputminted{scala}{\code/strategy/TextFormatter.kt}

The above interface contains only one method, \verb!format!, used to format the text.

Some sample test code might look like:

\inputminted{scala}{\code/strategy/TestStrategyPattern.kt}

The above code will result to the following output:

\begin{Verbatim}
[CapTextFormatter]: TESTING TEXT IN CAPS FORMATTER
[LowerTextFormatter]: testing text in lower formatter	
\end{Verbatim}

%\begin{solution}
%\inputminted{scala}{\codesol/strategy/CapTextFormatter.kt}
%
%The above class, \verb!CapTextFormatter!, is a concrete text formatter 
%that implements the \verb!TextFormatter! interface and the
%class is used to change the text into capital case.	
%
%\inputminted{scala}{\codesol/strategy/LowerTextFormatter.kt}
%
%The \verb!LowerTextFormatter! is a concrete text formatter that implements the 
%\verb!TextFormatter !interface and the class is used to change the text into small case.
%
%\inputminted{scala}{\codesol/strategy/TextEditor.kt}
%
%The above class is the \verb!TextEditor! class which holds a reference to the 
%\verb!TextFormatter! interface. The class contains the \verb!publishText! 
%method which forwards the text to the formatter to publish the text in the desired format.
%\end{solution}

\question %Abstract Factory Design Pattern
A product company, \textsc{Bigfish}, has changed the way they take orders from their clients. 
The company uses an application to take orders from them. 
They receive orders, errors in orders, feedback for the previous order, 
and responses to the order, in an XML format. 

Now the clients donâ't want to follow the companyâ's specific XML rules. 
The clients want to use their own XML rules to communicate with \textsc{Bigfish}. 
This means that for every client, the company should have client specific XML parsers. 
For example, for the NYC client there should be four specific types of XML parsers, i.e.,
\verb!NYCErrorXMLParser!, \verb!NYCFeedbackXML!, \verb!NYCOrderXMLParser!, 
\verb!NYCResponseXMLParser!, and four different parsers for the London client.

The company has asked you to change the application according to the new requirements. 
To develop the parser for the original application they used a \textsc{Factory 
Method} design pattern in which the exact object to use is decided by the subclasses 
according to the type of parser. Now, to implement the new requirements, 
you are required to use a factory of factories, i.e., an \textsc{Abstract Factory}.

\textbf{Note:} You will need parsers according to client specific XMLs, so you will create 
different factories for different clients which will provide the client specific XML to parse. 
You will achieve this by creating an \textsc{Abstract Factory} and then implement the 
factory to provide a client specific XML factory. 
Then you will use that factory to get the desired client specific XML parser object.

To implement the pattern you first need to create an interface that will be 
implemented by all the concrete factories:

\inputminted{scala}{\code/abstractfactory/AbstractParserFactory.kt}

The above interface is implemented by the client specific concrete factories which will 
provide the XML parser object to the client object. The \verb!getParserInstance! method 
takes the \verb!parserType !as an argument which is used to get the message specific 
(error parser, order parser etc) parser object.

The following is a test class for the resulting code:
\inputminted{scala}{\code/abstractfactory/TestAbstractFactoryPattern.kt}
The above code will result to the following output:
\begin{Verbatim}
NYC Parsing order XML...
NYC Order XML Message
************************************
London Parsing feedback XML...
London Feedback XML Message
\end{Verbatim}

%\begin{solution}
%\inputminted{scala}{\codesol/abstractfactory/NYCParserFactory.kt}
%
%\inputminted{scala}{\codesol/abstractfactory/LondonParserFactory.kt}
%	
%\inputminted{scala}{\codesol/abstractfactory/XMLParser.kt}
%
%\inputminted{scala}{\codesol/abstractfactory/NYCErrorXMLParser.kt}
%
%\inputminted{scala}{\codesol/abstractfactory/NYCFeedbackXMLParser.kt}
%
%\inputminted{scala}{\codesol/abstractfactory/NYCOrderXMLParser.kt}
%
%\inputminted{scala}{\codesol/abstractfactory/NYCResponseXMLParser.kt}
%
%\inputminted{scala}{\codesol/abstractfactory/LondonErrorXMLParser.kt}
%
%\inputminted{scala}{\codesol/abstractfactory/LondonFeedbackXMLParser.kt}
%
%\inputminted{scala}{\codesol/abstractfactory/LondonOrderXMLParser.kt}
%
%\inputminted{scala}{\codesol/abstractfactory/LondonResponseXMLParser.kt}
%
%To avoid a dependency between the client code and the factories, optionally we can implement
%a \emph{factory-producer} which has a \verb!static! method and is responsible for 
%providing a desired factory object to the client object.
%
%\inputminted{scala}{\codesol/abstractfactory/ParserFactoryProducer.kt}
%\end{solution}

\question
As an example of the \textsc{Builder Pattern}, consider a Car company which displays its 
different cars to its customers using a graphical model. 
The company has a graphical tool which displays the car on the screen. 
The requirements of the tool are that it is provided with a car object display. 
The car object should contain the car's specifications. 
The graphical tool uses these specifications to display the car. 

The company has classified its cars into different groups, e.g., Sedan, or Sports Car. 
There is only one car object, and our job is to create the car object according 
to the classification. For example, for a sedan car, a car object 
according to the sedan specification should be built or, if a sports car is required, 
then a car object according to the sports car specification should be built. 

Currently, the Company wants only these two types of cars, 
but it may require other types of cars also in the future.
You are required to create two different builders, one of each classification, 
i.e., for sedan and sports cars. 
The two builders will help in building the car object according to its specification. 

Below is the \verb!Car! class which contains some of the important components 
of the car that are required to construct the complete car object.

\inputminted{scala}{\code/builder/Car.kt}

You should be able to test your implementation using the following 
\verb!TestBuilderPattern! class:

\inputminted{scala}{\code/builder/TestBuilderPattern.kt}

The above code will result to the following output:
\begin{Verbatim}
--------------SEDAN--------------------- 
 Body: External dimensions: overall length (inches): 202.9, 
 	overall width (inches): 76.2, overall height (inches): 60.7, 
 	wheelbase (inches): 112.9, front track (inches): 65.3, 
 	rear track (inches): 65.5 and curb to curb turning circle (feet): 39.5
 Power: 285 hp @ 6,500 rpm; 253 ft lb of torque @ 4,000 rpm
 Engine: 3.5L Duramax V 6 DOHC
 Breaks: Four-wheel disc brakes: two ventilated. Electronic brake distribution
 Seats: Front seat centre armrest.Rear seat centre armrest.Split-folding rear seats
 Windows: Laminated side windows.Fixed rear window with defroster
 Fuel Type: Diesel 19 MPG city, 29 MPG highway, 23 MPG combined and 437 mi. range
--------------SPORTS--------------------- 
 Body: External dimensions: overall length (inches): 192.3, 
 	overall width (inches): 75.5, overall height (inches): 54.2, 
 	wheelbase (inches): 112.3, front track (inches): 63.7, 
 	rear track (inches): 64.1 and curb to curb turning circle (feet): 37.7
 Power: 323 hp @ 6,800 rpm; 278 ft lb of torque @ 4,800 rpm
 Engine: 3.6L V 6 DOHC and variable valve timing
 Breaks: Four-wheel disc brakes: two ventilated. Electronic brake distribution. 
 	Stability control
 Seats: Driver sports front seat with one power adjustments manual height, 
 	front passenger seat sports front seat with one power adjustments
 Windows: Front windows with one-touch on two windows
 Fuel Type: Petrol 17 MPG city, 28 MPG highway, 20 MPG combined and 380 mi. range	
\end{Verbatim}

%\begin{solution}
%The \verb!CarBuilder! is the builder interface and contains a set of common methods 
%used to build the car object and its components.	
%
%\inputminted{scala}{\codesol/builder/CarBuilder.kt}
%
%The \verb!getCar! method is used to return the final car object to the client 
%after its construction.
%
%Two implementations of the \verb!CarBuilder! interface, one for each type of car, 
%i.e., for sedan and sports car.
%
%\inputminted{scala}{\codesol/builder/SedanCarBuilder.kt}
%\inputminted{scala}{\codesol/builder/SportsCarBuilder.kt}
%\end{solution}

\question % Facade design pattern
\textsc{BetterBisc} is product based company and it has launched a product in the market, 
named \emph{Schedule Server}. It is a kind of server in itself, and it is used to manage jobs. 
The jobs could be any kind of jobs like sending a list of emails, sms, reading or writing 
files from a destination, or just simply transferring files from a source to the destination. 

The product is used by the developers to manage such types of job and enables them to concentrate 
on their business goal. The server executes each job at their specified time and 
also manages any underling issues like concurrency, security, etc. As a developer, one just needs
to code the relevant business requirements and a good amount of API calls is 
provided to schedule a job according to their needs.

Everything was going fine, until the clients started complaining about starting and stopping the
the server. Although the server is working well, the initialisation and the shut down 
processes are very complex and they want an easy way to do that. 
The server has exposed a complex interface to the clients which looks a bit \emph{busy} 
to them.

We need to provide an easy way to start and stop the server.

A complex interface to the client is already considered as a fault in the design of the 
current system.  Fortunately or unfortunately, we cannot start the design and the coding
from scratch. We need a way to resolve this problem and make the interface easy to access.

The problem faced by the clients in using the schedule server is 
the complexity brought by the server in order to start and stop its services. 
The client wants a simple way to do it. The following is the code that clients 
required to write to start and stop the server.

\begin{minted}{scala}
ScheduleServer scheduleServer = new ScheduleServer();	
\end{minted}
To start the server, the client needs to create an object of the 
\mintinline{scala}{ScheduleServer} class and then needs to call the following methods (in 
sequence) to start and initialise the server.
\begin{minted}{scala}
scheduleServer.startBooting();
scheduleServer.readSystemConfigFile();
scheduleServer.init();
scheduleServer.initializeContext();
scheduleServer.initializeListeners();
scheduleServer.createSystemObjects();
System.out.println("Start working......");
System.out.println("After work done.........");	
\end{minted}
To stop the server, the client needs to call the following methods in sequence:
\begin{minted}{scala}
scheduleServer.releaseProcesses();
scheduleServer.destory();
scheduleServer.destroySystemObjects();
scheduleServer.destoryListeners();
scheduleServer.destoryContext();
scheduleServer.shutdown();	
\end{minted}
To resolve this, you will create a \emph{facade class} which will \emph{wrap} a server object. 
This class will provide simple interfaces (methods) for the client. 
These interfaces internally will call the methods on the server object. 

\inputminted{scala}{\code/facadepattern/ScheduleServer.kt}

%\begin{solution}
%\inputminted{scala}{\codesol/facadepattern/ScheduleServerFacade.kt}	
%\end{solution}


\question % Bridge Design Pattern
\textsc{Dodgy Security Systems} is a security and electronic company which produces and 
assembles products for cars. It delivers any car electronic or security system you want, 
from air bags to GPS tracking system, reverse parking system, etc. 
The company uses a well defined object oriented approach to keep track of their products 
using software which is developed and maintained only by them. 
They receive a vehicle, produce the system for it, and then assemble it and install it into 
the vehicle.

Recently, they received new orders from \textsc{BigWheel} (a car company) to produce 
central locking and gear lock systems for their new car model. To maintain this, they are 
creating a new software system. They started by creating a new abstract class 
\verb!CarProductSecurity!, in which they kept some car specific methods and some of the 
features which they thought are common to all security products. 
Then they extended the class and created two different sub classes:
\begin{itemize}
\item \verb!BigWheelCentralLocking!, and 
\item \verb!BigWheelGearLocking!. 
\end{itemize}
The class diagram is as follows:

\includegraphics[width=\textwidth]{images/design}

After a while, another car company \textsc{Motoren} asked them to produce a new system of 
central locking and gear lock for their new model. 
Since, the same security system cannot be used in both models of different cars,  
\textsc{Dodgy Security Systems} has produced a new system for them, and also has created 
two new classes 
\begin{itemize}
\item \verb!MotorenCentralLocking!, and 
\item \verb!MotorenGearLocking! 
\end{itemize}
which also extend the \verb!CarProductSecurity! class.
Now the new class diagram looks like this:

\includegraphics[width=\textwidth]{images/newdesign}

what happens if another car company demands another new system for central locking and gear lock? 
One would need to create another two new classes for it. 
This design will create one class per system, or worse, if a reverse parking system 
is produced for each of these two car companies, 
two more new classes will be created for each of them.

Provide a solution to this problem using an implementation of the \textsc{Bridge Design Pattern}.
We provide an appropriate test method:

\inputminted{scala}{\code/bridge/TestBridgePattern.kt}

which will produce the following output:

\begin{Verbatim}
Producing Central Locking System
Modifying product Central Locking System according to BigWheel xz model
Assembling Central Locking System for BigWheel xz model
Car: BigWheel xz model, Product:Central Locking System

Producing Gear Locking System
Modifying product Gear Locking System according to BigWheel xz model
Assembling Gear Locking System for BigWheel xz model
Car: BigWheel xz model, Product:Gear Locking System

Producing Central Locking System
Modifying product Central Locking System according to Motoren lm model
Assembling Central Locking System for Motoren lm model
Car: Motoren lm model, Product:Central Locking System

Producing Gear Locking System
Modifying product Gear Locking System according to Motoren lm model
Assembling Gear Locking System for Motoren lm model
Car: Motoren lm model, Product:Gear Locking System	
\end{Verbatim}

%\begin{solution}
%Instead of creating a subclass for each product per car model in the above discussed problem, 
%we can separate the design into two different hierarchies. One interface is for the 
%product which will be used as an implementer and the other will be an abstraction of car type. 
%The implementer will be implemented by the concrete implementers and provides an 
%implementation for it. On the other side, the abstraction will be extended by more 
%refined abstraction.	
%
%\includegraphics[width=\textwidth]{images/reviseddesign}
%
%\inputminted{scala}{\codesol/bridge/Product.kt}
%
%The implementer \verb!Product! has a method \verb!produce()! which will be used by the 
%concrete implementers to provide concrete functionality to it. The method will 
%produce the base model of the product which can be used with any 
%car model after some modifications specific to that car model.
%
%\inputminted{scala}{\codesol/bridge/CentralLocking.kt}
%\inputminted{scala}{\codesol/bridge/GearLocking.kt}
%
%The two different concrete implementers provide implementation to the \verb!Product! implementer.
%
%Now the abstraction, the \verb!Car! class which holds a reference of a product type 
%and provides two abstract methods \verb!produceProduct()! and \verb!assemble()!.
%
%\inputminted{scala}{\codesol/bridge/Car.kt}
%
%The subclasses of the \verb!Car! will provide the concrete and specific implementation 
%to the methods \verb!assemble()! and produce \verb!Product()!.
%
%\inputminted{scala}{\codesol/bridge/BigWheel.kt}
%\inputminted{scala}{\codesol/bridge/Motoren.kt}
%\end{solution}

\question
HTML is hierarchical in nature, its starts from an \verb!<html>! tag which is the 
parent or the root tag, and it contains other tags which can be a parent or a child tag.
The \textsc{Composite Pattern} in Java can be implemented using the component class as an 
abstract class or an interface. In this question, you will use an abstract class which 
contains all the important methods used in a composite class and a leaf class.

\inputminted{scala}{\code/composite/HtmlTag.kt}

Given the above class complete two subclasses of the class:
\begin{description}
	\item[HtmlParentElement] --- handles the \emph{child} nodes, and
	\item[HtmlElement] --- handles the \emph{leaf} nodes.
\end{description}
The following class should be used to test out your classes.

\inputminted{scala}{\code/composite/TestCompositePattern.kt}

The above code will result to the following output:
\begin{Verbatim}
<html>
<body>
<P>Testing html tag library</P>
<P>Paragraph 2</P>
</body>
</html>	
\end{Verbatim}

%\begin{solution}
%\inputminted{scala}{\codesol/composite/HtmlElement.kt}
%\inputminted{scala}{\codesol/composite/HtmlParentElement.kt}	
%\end{solution}
\end{questions}

\end{document}